\subsection{Protocol-locked confirmation on two new sources}\label{sec:newdesign}
After the preceding sensitivities, we fixed an additional protocol before numeric data
acquisition or calibration inspection. It selects MPI-BGC Jena roof temperature and
pressure for the first 3,600 hours of 2026, and PM2.5 at the Aotizhongxin and Dingling
stations for the first 3,600 hours from 1 March 2013 in the UCI Beijing archive
\citep{jena2026weather,chen2017air}. These are four series from two sources, evaluated
with the same two cached Chronos checkpoints: eight dependent model--series pairs.
Both models emit all 24 horizons at 2,976 consecutive origins. The main test segment
has 1,776 origins before excluding missing targets. Context length, batch size, dtype,
checkpoint-file hashes, timestamps and generation code are recorded for every archive.

Missing input observations are forward filled causally; missing targets are excluded
from score memory, controller increments and evaluation. Windows contain the latest
$W$ observed, revealed scores. Issued-state phases continue without an increment when
their outcome is missing. Local coverage uses 150 calendar-origin windows with their
observed denominators. The weather forecast targets are complete; the one missing
initial weather hour precedes all forecast contexts. The new-source protocol is locally
timestamped and hashed, not externally preregistered. Public historical air-quality data
may have been used in neural pretraining; new-to-calibration does not imply unseen-to-model.

For both scores and states, nested grids comprise the original rates at $W=450$;
the wide grid described below; and that grid augmented with zero
and four smaller rates ($10^{-6}$, $5\times10^{-6}$, $10^{-5}$ and $5\times10^{-5}$).
The final stage crosses these expanded rates with $W\in\{200,450,900\}$. Validation selects the minimum interval
score, breaking ties toward smaller rate and then window. Zero-rate trajectories
provide both fixed-window and selected-window rolling controls. Early and late
validation intervals, $[600+h-1,900-h+1)$ and $[900,1200-h+1)$, probe selector
stability without changing the primary full-validation selector or test segment.

The wide grid combines rates 0.0001, 0.0002, 0.0005, 0.001, 0.002, 0.005,
0.01, 0.02, 0.05 and 0.1 with the three original rates divided by $h$.

The predefined complete-weather subset also receives full AcMCP scorecasting: the
pinned author's MA$(h-1)$ plus cross-horizon regression implementation, at all horizons
1--24, with its supported CSS fitting option. We identify it as AcMCP-CSS, not the
default CSS-ML configuration. Rolling history is 600, $C_{\rm sat}=0.5$, and $K_I$
uses only the first 600 revealed target-clock error rows. The rate grid is
$\{0.001,0.005,0.01,0.05,0.1\}$ and these rates divided by $h$; the same validation
block selects it. Fitted scorecasts are independent of this rate, so they are cached
once per archive. A cached interval replay must agree with every original R trajectory
at rate 0.01 before evaluating the other rates. A matched no-scorecast PID control
uses the identical rate grid. Both are signed-error complete configurations; their
comparison to native CQR does not isolate controller quality.

\subsection{PV support and asymmetric-residual controls}\label{sec:pvcontrol}
On all 32 original one-step PV pairs, we additionally compare the established
signed-error construction
$[m_t+Q_{a_t/2}(Y-m),m_t+Q_{1-a_t/2}(Y-m)]$ with symmetric residuals and CQR.
It shares the total-miscoverage ACI state, timing and window. Each construction
uses the original nine rate/window settings; fixed defaults and a separately selected
zero-rate rolling control are retained. This is exploratory analysis of previously
inspected data, distinct from the new-source confirmation.

All three constructions are evaluated raw and after the same known-support projection
$[L,U]\mapsto[\max(0,L),\max(0,U)]$. Projected and raw configurations are separately
selected on validation score, while each candidate keeps its raw feedback trajectory:
the projection isolates interval geometry and is not a new feedback algorithm.
For a nonnegative target and finite ordered endpoints, projection cannot worsen
interval score or lose a covered target. In particular, when $L<0\leq U$, score
improvement is exactly $-L$, with unchanged coverage. When $U<0$, both endpoints
become zero and previously uncovered zeros can become covered. We therefore report
coverage changes as well as width and miss-distance penalties.

In addition to actual zero/positive targets, an observable grouping compares each point
forecast with the median of the first 600 point forecasts. This threshold is fixed
before validation; the high-forecast group is known at issue time and is not selected
using the future outcome. It is a diagnostic control, not a conditional-validity guarantee.
